La stima delle matrici di correlazione fra asset finanziari a partire dalle sole serie storiche presenta due limiti strutturali: il rumore statistico che degrada le matrici empiriche quando il numero di titoli è dello stesso ordine di grandezza della lunghezza della serie, e il fatto che il passato offra una sola traiettoria fra tutte quelle possibili. Questa tesi, svolta in collaborazione con Intesa Sanpaolo, affronta entrambi i limiti testando l'efficacia dei Variational Autoencoder (VAE) come modello generativo stocastico per la simulazione di scenari di mercato realistici.

Il lavoro parte dai rendimenti giornalieri di 362 titoli dell'indice S\&P 500 quotati con continuità nel ventennio 2000-2020. Da essi vengono estratte, con finestra mobile di 724 giorni e passo di 10, 461 matrici di correlazione empiriche, ridotte a 412 da un filtro di validità basato sulla proprietà di Perron-Frobenius. Poiché una matrice prodotta da un modello generativo non è automaticamente semidefinita positiva --- vincolo imprescindibile per una matrice di correlazione valida --- l'input dei modelli non è la matrice grezza, bensì il fattore triangolare inferiore della sua decomposizione di Cholesky: qualunque vettore prodotto dal decoder si ricompone così in una matrice semidefinita positiva per costruzione.

La metodologia segue un percorso incrementale su quattro architetture confrontate a parità di dimensione latente: Principal Component Analysis come benchmark lineare, Autoencoder Lineare, Autoencoder Profondo e Variational Autoencoder. La valutazione combina metriche di errore spaziale (MSE, MAE, norma di Frobenius) e metriche topologiche costruite sul Minimum Spanning Tree. Sul piano della pura ricostruzione storica l'Autoencoder Profondo domina tutti i modelli a ogni dimensione latente; il VAE, addestrato con una loss di ricostruzione basata sulla Divergenza di Jeffreys fra le distribuzioni associate alle matrici, si colloca in posizione intermedia: a venti dimensioni latenti supera sia la PCA sia l'Autoencoder Lineare, pur restando dietro all'Autoencoder deterministico. La regolarizzazione variazionale non compromette dunque in modo sostanziale la fedeltà di ricostruzione.

È sul piano generativo che il VAE mostra il proprio valore distintivo: la continuità e la densità dello spazio latente indotte dalla Divergenza di Kullback-Leibler sono precluse a un autoencoder deterministico, il cui spazio latente resta frammentario. La procedura adottata non è un semplice campionamento dal prior, bensì un block-bootstrap storico degli incrementi latenti giornalieri, che ancora ciascuno scenario allo stato di mercato osservato oggi e produce una distribuzione di evoluzioni plausibili della struttura di correlazione a un orizzonte di un mese, valutabile in ottica di rischiosità.

I mille scenari generati sono stati validati rispetto ai cinque fatti stilizzati delle matrici di correlazione finanziarie: distribuzione dei coefficienti spostata verso valori positivi, spettro degli autovalori conforme alla legge di Marchenko-Pastur, proprietà di Perron-Frobenius, struttura gerarchica settoriale e topologia scale-free del Minimum Spanning Tree. Tutti e mille gli scenari soddisfano simultaneamente i cinque criteri, con similarità coseno superiore a 0.98 rispetto alla matrice odierna sui modi principali. L'architettura proposta si conferma quindi uno strumento autonomo per la generazione di ecosistemi finanziari sintetici e matematicamente coerenti.
